\section{Majuskuły}
\label{sec:majuskuy}

Majuskuły są wypisane w pierwszym wierszu zestawu.

\begin{enumerate}
\item Litera `A --- patrz przykładowe wystąpienia na
  il. \vref{fig:Ungler1_pismo01_A}.
  % W standardzie Unicode znak
  % \ucode{0041} \uname{LATIN CAPITAL LETTER A}.


\begin{figure}[ht]
  \centering
  \includegraphics[height=10ex]{img/A1511p14A}
  \includegraphics[height=10ex]{img/Additio1511p8b}
  \includegraphics[height=10ex]{img/Articulus1511p8a}
  \caption{`A: \Ju{A}, \Ju{Additio}, \Ju{Articulus} (\textit{Algorithmus})}
  \label{fig:Ungler1_pismo01_A}
\end{figure}

\item Litera `B --- patrz jedyne znalezione  wystąpienie na
  il. \vref{fig:Ungler1_pismo01_B}.
  % W standardzie Unicode znak
  % \ucode{0042} \uname{LATIN CAPITAL LETTER B}.


\begin{figure}[ht]
  \centering
  \includegraphics[height=10ex]{img/A511p20Boetius}
  \caption{`B: \Ju{Boetius} (\textit{Algorithmus})}
  \label{fig:Ungler1_pismo01_B}
\end{figure}

\item Litera `C --- patrz przykładowe wystąpienia na
  il. \vref{fig:Ungler1_pismo01_C}.
  % W standardzie Unicode znak
  % \ucode{0043} \uname{LATIN CAPITAL LETTER C}.

  \begin{figure}[ht]
  \centering
  \includegraphics[height=10ex]{img/A1521p9Cum1}
  \includegraphics[height=10ex]{img/A1521p10Centum}
%  \includegraphics[height=10ex]{img/A1521p20Cum3}
  \includegraphics[height=10ex]{img/A1521p8Compositus}
  \caption{`C: \Ju{Cum}, \Ju{Centum}, \Ju{Compositus} (\textit{Algorithmus})}
    \label{fig:Ungler1_pismo01_C}
\end{figure}

\item Litera `D --- patrz przykładowe wystąpienia na
  il. \vref{fig:Ungler1_pismo01_D}.
  % W standardzie Unicode znak
  % \ucode{0044} \uname{LATIN CAPITAL LETTER D}.


\begin{figure}[ht]
  \centering
  \includegraphics[height=10ex]{img/A1521p21Deinde}
  \includegraphics[height=10ex]{img/A1521p8Digitus}
  \includegraphics[height=10ex]{img/A1521p9Decima}
  \includegraphics[height=10ex]{img/A1521p9Differentia}
  \caption{`D: \Ju{Deinde}, \Ju{Digitus}, \Ju{Decima},
    \Ju{Differrntia}[sic] (\textit{Algorithmus})}
  \label{fig:Ungler1_pismo01_D}
\end{figure}



\item Litera `E --- patrz przykładowe wystąpienia na
  il. \vref{fig:Ungler1_pismo01_E}.
  % W standardzie Unicode znak
  % \ucode{0045} \uname{LATIN CAPITAL LETTER E}.


\begin{figure}[ht]
  \centering
  \includegraphics[height=10ex]{img/A1521p8Est}
  \includegraphics[height=10ex]{img/A1521p13Et}
  \includegraphics[height=10ex]{img/A1521p14Ex}
%  \includegraphics[height=10ex]{img/A1521p9Differentia}
  \caption{`E: \Ju{Est}, \Ju{Et}, \Ju{Ex}
    (\textit{Algorithmus})}
  \label{fig:Ungler1_pismo01_E}
\end{figure}




\item Litera `F --- nie znalazłem żadnego wystąpienia tego znaku.

\item Litera `G --- nie znalazłem żadnego wystąpienia tego znaku.

\item Litera `H --- patrz przykładowe wystąpienia na
  il. \vref{fig:Ungler1_pismo01_H}.
  % W standardzie Unicode znak
  % \ucode{0048} \uname{LATIN CAPITAL LETTER H}.


\begin{figure}[ht]
  \centering
  \includegraphics[height=10ex]{img/A1521p11Hoc}
  \includegraphics[height=10ex]{img/A1521p8Huius}
  \includegraphics[height=10ex]{img/A1521p17His}
  \includegraphics[height=10ex]{img/A1521p8Hanc}
   \caption{`H: \Ju{Hoc}, \Ju{Huius}, \Ju{His} (?), \Ju{Hanc}
 %  (\textit{Algorithmus})}
   (Algorithmus)}
  \label{fig:Ungler1_pismo01_H}
\end{figure}

\item Litera `I --- patrz przykładowe wystąpienia na
  il. \vref{fig:Ungler1_pismo01_I}.
  % W standardzie Unicode znak
  % \ucode{0049} \uname{LATIN CAPITAL LETTER I}.

\begin{figure}[ht]
  \centering
  \includegraphics[height=10ex]{img/A1521p10In}
  \includegraphics[height=10ex]{img/A1521p11Iste}
  \includegraphics[height=10ex]{img/A1521p10Igitur}
  \includegraphics[height=10ex]{img/A1521p25Impressus}
  \caption{`I: \Ju{In}, \Ju{Iste}, \Ju{Igitur}, \Ju{Impressus}
    (\textit{Algorithmus})}
  \label{fig:Ungler1_pismo01_I}
\end{figure}

\item Litera `K --- nie znalazłem żadnego wystąpienia tego znaku.

\item Litera `L --- patrz jedyne znalezione  wystąpienie na
  il. \vref{fig:Ungler1_pismo01_L}.
  % W standardzie Unicode znak
  % \ucode{004C} \uname{LATIN CAPITAL LETTER L}.


\begin{figure}[ht]
  \centering
  \includegraphics[height=10ex]{img/A1521p9Locus}
  % \includegraphics[height=10ex]{img/A1521p11Iste}
  % \includegraphics[height=10ex]{img/A1521p10Igitur}
  % \includegraphics[height=10ex]{img/A1521p25Impressus}
  \caption{`L: \Ju{Locus}
%    , \Ju{Iste}, \Ju{Igitur}, \Ju{Impressus}
    (\textit{Algorithmus})}
  \label{fig:Ungler1_pismo01_L}
\end{figure}

\item Litera `M --- patrz przykładowe wystąpienia na
  il. \vref{fig:Ungler1_pismo01_M}.
  % W standardzie Unicode znak
  % \ucode{004D} \uname{LATIN CAPITAL LETTER M}.


\begin{figure}[ht]
  \centering
  \includegraphics[height=10ex]{img/A1521p20Maius}
  \includegraphics[height=10ex]{img/A1521p20Minus}
  \includegraphics[height=10ex]{img/A1521p8Mediatio}
  \includegraphics[height=10ex]{img/A1521p8Materialiter}
  \caption{`M: \Ju{Maius}, \Ju{Minus}, \Ju{Mediatio}, \Ju{Materialiter}
    (\textit{Algorithmus})}
  \label{fig:Ungler1_pismo01_M}
\end{figure}

\item Litera `N --- patrz przykładowe wystąpienia na
  il. \vref{fig:Ungler1_pismo01_N}.
  % W standardzie Unicode znak
  % \ucode{004E} \uname{LATIN CAPITAL LETTER N}.


\begin{figure}[ht]
  \centering
  \includegraphics[height=10ex]{img/A1521p21Nec}
  \includegraphics[height=10ex]{img/A1521p10Numerum}
  \includegraphics[height=10ex]{img/A1521p9Notandum}
  \includegraphics[height=10ex]{img/A1521p8Numeratio}
  \caption{`N: \Ju{Nec}, \Ju{Numerum}, \Ju{Notandum}, \Ju{Numeratio}
    (\textit{Algorithmus})}
  \label{fig:Ungler1_pismo01_N}
\end{figure}

\item Litera `O --- patrz przykładowe wystąpienia na
  il. \vref{fig:Ungler1_pismo01_O}.
  % W standardzie Unicode znak
  % \ucode{004F} \uname{LATIN CAPITAL LETTER O}.


\begin{figure}[ht]
  \centering
  \includegraphics[height=10ex]{img/A1521p9Omnis2}
  \includegraphics[height=10ex]{img/A1521p9Omnis1}
  % \includegraphics[height=10ex]{img/A1521p9Notandum}
  % \includegraphics[height=10ex]{img/A1521p8Numeratio}
  \caption{`O: \Ju{Omnis}
    % , \Ju{Numerum}, \Ju{Notandum}, \Ju{Numeratio}
   (\textit{Algorithmus})}
  \label{fig:Ungler1_pismo01_O}
\end{figure}

\item Litera `P --- patrz przykładowe wystąpienia na
  il. \vref{fig:Ungler1_pismo01_P}.
  % W standardzie Unicode znak
  % \ucode{0050} \uname{LATIN CAPITAL LETTER P}.


\begin{figure}[ht]
  \centering
  \includegraphics[height=10ex]{img/A1521p8Progressio}
  \includegraphics[height=10ex]{img/A1521p8Primo}
 % \includegraphics[height=10ex]{img/A1521p15Prima}
  % \includegraphics[height=10ex]{img/A1521p8Numeratio}
  \caption{`P: \Ju{Progressio}, \Ju{Primo}
%    , \Ju{Prima}
   (\textit{Algorithmus})}
  \label{fig:Ungler1_pismo01_P}
\end{figure}

\item Litera `Q --- patrz przykładowe wystąpienia na
  il. \vref{fig:Ungler1_pismo01_Q} i na
    il. \vref{fig:Ungler1_pismo01_mmacron} (\Ju{Q??}).
  % W standardzie Unicode znak
  % \ucode{0051} \uname{LATIN CAPITAL LETTER Q}.


\begin{figure}[ht]
  \centering
  \includegraphics[height=10ex]{img/A1521p14Quarta}
  \includegraphics[height=10ex]{img/A1521p14Quinta}
  \includegraphics[height=10ex]{img/A1521p14Quinam}
  % \includegraphics[height=10ex]{img/A1521p8Numeratio}
  \caption{`Q: \Ju{Quarta},
  \Ju{Quinta}, \Ju{Quinam}?
   (\textit{Algorithmus})}
  \label{fig:Ungler1_pismo01_Q}
\end{figure}

\item Litera `R --- patrz przykładowe wystąpienia na
  il. \vref{fig:Ungler1_pismo01_R}.
  % W standardzie Unicode znak
  % \ucode{0052} \uname{LATIN CAPITAL LETTER R}.


\begin{figure}[ht]
  \centering
  \includegraphics[height=10ex]{img/A1521p22Radix}
  \includegraphics[height=10ex]{img/A1521p22Radicem}
  % \includegraphics[height=10ex]{img/A1521p9Notandum}
  % \includegraphics[height=10ex]{img/A1521p8Numeratio}
  \caption{`R: \Ju{Radix}, \Ju{Radicem}
   (\textit{Algorithmus})}
  \label{fig:Ungler1_pismo01_R}
\end{figure}

\item Litera `S --- patrz przykładowe wystąpienia na
  il. \vref{fig:Ungler1_pismo01_S} i na
    il. \vref{fig:Ungler1_pismo01_dd} (\Ju{Secundo}).
  % W standardzie Unicode znak
  % \ucode{0053} \uname{LATIN CAPITAL LETTER S}.


\begin{figure}[ht]
  \centering
  \includegraphics[height=10ex]{img/A1521p10Sexto}
  \includegraphics[height=10ex]{img/A1521p10Sinistrorsum}
  \includegraphics[height=10ex]{img/A1521p8Subtracio}
  \includegraphics[height=10ex]{img/A1511p1Similiter}
  \caption{`S: \Ju{Sexto}, \Ju{Sinistrorsum}, \Ju{Subtratio} [sic!],
\Ju{Similiter} (\textit{Algorithmus})}
  \label{fig:Ungler1_pismo01_S}
\end{figure}

\item Litera `T --- patrz przykładowe wystąpienia na
  il. \vref{fig:Ungler1_pismo01_T}.
  % W standardzie Unicode znak
  % \ucode{0054} \uname{LATIN CAPITAL LETTER T}.


\begin{figure}[ht]
  \centering
  \includegraphics[height=10ex]{img/A1521p10Tres}
  \includegraphics[height=10ex]{img/A1521p15Tercia}
  \includegraphics[height=10ex]{img/A1521p8Tertio}
%  \includegraphics[height=10ex]{img/A1511p1Similiter}
  \caption{`T: \Ju{Tres}?, \Ju{Tercia}, \Ju{Tertio}
    (\textit{Algorithmus})}
  \label{fig:Ungler1_pismo01_T}
\end{figure}

\item Litera `U --- patrz przykładowe wystąpienia na
  il. \vref{fig:Ungler1_pismo01_U}.
  % W standardzie Unicode znak
  % \ucode{0055} \uname{LATIN CAPITAL LETTER U}.


\begin{figure}[ht]
  \centering
  \includegraphics[height=10ex]{img/A1521p8Unde}
  \includegraphics[height=10ex]{img/A1521p8Unitas}
  \includegraphics[height=10ex]{img/A1521p15Uel}
  \includegraphics[height=10ex]{img/A1521p19Un_}
  \caption{`U: \Ju{Unde}, \Ju{Unitas}, \Ju{Uel}, \Ju{Unde}?
    (\textit{Algorithmus})}
  \label{fig:Ungler1_pismo01_U}
\end{figure}

\item Ostatnia majuskuła to ,,ogoniaste Z'' (ang. \textit{tailed Z});
  nie znalazłem jej wystąpienia w tekście. W standardzie Unicode od
  wersji 3.0 jest ona reprezentowana przez \ucode {01B7} \uname{LATIN
    CAPITAL LETTER EZH} (\Ju{Ʒ}) \autocite[113-114]{haugen15:_mufi1}.

\end{enumerate}


\section{Minuskuły}
\label{sec:minuskuy}

\begin{enumerate}
\item Wiersz pierwszy (w wykazie drugi):
\begin{enumerate}
\item Litera `a --- patrz przykładowe wystąpienia na
  il. \vref{fig:Ungler1_pismo01_D} (\Ju{Decima}),
  il. \vref{fig:Ungler1_pismo01_H} (\Ju{Hanc}),
  il. \vref{fig:Ungler1_pismo01_M} (\Ju{Maius}, \Ju{Mediatio}, \Ju{Materialiter}),
  il. \vref{fig:Ungler1_pismo01_N} (\Ju{Notandum}, \Ju{Numeratio}),
  il. \vref{fig:Ungler1_pismo01_Q} (\Ju{Quarta}, \Ju{Quinta}),
  il. \vref{fig:Ungler1_pismo01_R} (\Ju{Radix}, \Ju{Radicem}),
  il. \vref{fig:Ungler1_pismo01_S} (\Ju{Subtratio} [sic!]),
  il. \vref{fig:Ungler1_pismo01_T} (\Ju{Tercia}),
  il. \vref{fig:Ungler1_pismo01_U} (\Ju{Unitas}),
  il. \vref{fig:Ungler1_pismo01_emacron} (\Ju{pertactionem}),
  il. \vref{fig:Ungler1_pismo01_f} (\Ju{figuras}, \Ju{formata}),
    il. \vref{fig:Ungler1_pismo01_h} (\Ju{habent}, \Ju{theca}),
il. \vref{fig:Ungler1_pismo01_ij} (\Ju{[ne]cessarii}, \Ju{aliis}),
  il. \vref{fig:Ungler1_pismo01_con} (\Ju{considerata}) i
 il. \vref{fig:Ungler1_pismo01_rum} (\Ju{linearum}).


\item Litera \textit{a} z kreską `{ā} --- patrz przykładowe wystąpienia na
  il. \vref{fig:Ungler1_pismo01_amacron}; kreskę nad literą
  interpretuję jako makron, choć w każdym przykładzie wygląda nieco
  inaczej (w standardzie Unicode \Ju{◌̄} \ucode{0304} \uname{COMBINING
    MACRON}); oznacza ona skrócenie.
\begin{figure}[ht]
  \centering
  \includegraphics[height=10ex]{img/A1521p10quamlibet}
  \includegraphics[height=10ex]{img/A1521p9illam}
  \includegraphics[height=10ex]{img/A1521p8quidam}
%  \includegraphics[height=10ex]{img/A1521p19Un_}
  \caption{`{ā}: \Ju{quamlibet}, \Ju{illam}, \Ju{quidam}
    (\textit{Algorithmus})}
  \label{fig:Ungler1_pismo01_amacron}
\end{figure}

\item Litera `b --- patrz przykładowe wystąpienie na
  il. \vref{fig:Ungler1_pismo01_S} (\Ju{Subtratio} [sic!]),
  il. \vref{fig:Ungler1_pismo01_h} (\Ju{habent}) i
    il. \vref{fig:Ungler1_pismo01_mmacron} (\Ju{omnibus}).

  \item Litera \textit{b} z trudnym do zidentyfikowania diakrytem
    (przecinek?) --- być może jednak jest to po prostu rozwidlenie
    laski, o którym pisze Porębski
    \autocite*[s. 11]{porębski05:_paleog} (Krzysztof Nowak zwrócił
    moją uwagę na książkę
    \autocite{derolez06:_palaeog_gothic_manus_books}, gdzie ten kształ
    jest opisany jako \textit{bifurcation}); znalazłem tylko jedno jej
    samodzielne wystąpienie w tekście, patrz
    il. \vref{fig:Ungler1_pismo01_bx}; w wydaniu z 1513 r. w tym samym
    słowie występuje zwykła litera \textit{b}.  Wiadomo (patrz np.
    tabele znaków na s. 53 i 62
    \autocite{perrousseaux05:_histoir_typog_guten_xviiie} i s. 32-33
    \autocite{schwenke00:_unter}), że litera ta występowała już w
    biblii Gutenberga, ale nie znam jej funkcji.  Wchodzi ona w skład
    omawianej dalej ligatury występującej w
    \autocite{sacro11:_algor_ioann_de_sacro_busto} wielokrotnie.
% Si autem pro
% bare velis vtrum

% Krzysztof Nowak <krzysztof.nowak@ijp.pan.pl>    
% Date: Sun, 21 Feb 2021 16:26:26 +0100
%     1. Chodzi chyba o bifurkację, "rozszczepienie", którego ilustrację
%     znaleźć może Pan Profesor np. na załączonym obrazku (4-5). Książka
%     A. Deroleza, The Palaeography of Gothic Manuscript Books: From the
%     Twelfth to the Early Sixteenth Century jest bardzo przydatna i
%     niezwykle szczegółowa. Dostępna jest zresztą w sieci w, nazwijmy
%     to, nie do końca legalnych źródłach
    

      % s. 11

      %           Prowadzi to do powstania pisma gotyckiego. Trzonki liter i, u
      %     oraz laski liter b, d, h, I u wierzchołków rozwidlają się,

\begin{figure}[ht]
  \centering
  \includegraphics[height=10ex]{img/A1521p12probare}
  % \includegraphics[height=10ex]{img/A1521p9illam}
  % \includegraphics[height=10ex]{img/A1521p8quidam}
%  \includegraphics[height=10ex]{img/A1521p19Un_}
  \caption{\Ju{probare}
    (\textit{Algorithmus})}
  \label{fig:Ungler1_pismo01_bx}
\end{figure}

\item Litera `c--- patrz przykładowe wystąpienia na
  il. \vref{fig:Ungler1_pismo01_A} (\Ju{Articulus}),
  il. \vref{fig:Ungler1_pismo01_D} (\Ju{Decima}),
  il. \vref{fig:Ungler1_pismo01_H} (\Ju{Hoc}, \Ju{Hanc}),
  il. \vref{fig:Ungler1_pismo01_L} (\Ju{Locus}),
  il. \vref{fig:Ungler1_pismo01_N} (\Ju{Nec}),
  il. \vref{fig:Ungler1_pismo01_R} (\Ju{Radicem}),
  il. \vref{fig:Ungler1_pismo01_T} (\Ju{Tercia}),
  il. \vref{fig:Ungler1_pismo01_emacron} (\Ju{pertactionem}),
    il. \vref{fig:Ungler1_pismo01_h} (\Ju{theca}),
 il. \ref{fig:Ungler1_pismo01_dd} (\Ju{Secundo}, \Ju{secundae}?),
   il. \vref{fig:Ungler1_pismo01_omacron} (\Ju{econuerso}),
il. \vref{fig:Ungler1_pismo01_p_} (\Ju{superficies}),
il. \vref{fig:Ungler1_pismo01_qet} (\Ju{quocienscumque)} i
 il. \vref{fig:Ungler1_pismo01_rum} (\Ju{locorum}).

\item Kolejna pozycja to ,,okragłe \textit{d}'', inaczej \textit{d}
  insularne (w standardzie Unicode \Ju{ꝺ} \ucode{A77A} \uname{LATIN
    SMALL LETTER INSULAR D}) --- patrz przykładowe wystąpienia na
  % il. \ref{fig:Ungler1_pismo01_A} (\Ju{Additio}),
  % il. \ref{fig:Ungler1_pismo01_D} (\Ju{Deinde}),
  % il. \ref{fig:Ungler1_pismo01_M} (\Ju{Mediatio}),
  % il. \ref{fig:Ungler1_pismo01_N} (\Ju{Notandum}),
  % il. \ref{fig:Ungler1_pismo01_R} (\Ju{Radix}, \Ju{Radicem}),
  % il. \ref{fig:Ungler1_pismo01_U} (\Ju{Undi}),
  % il. \ref{fig:Ungler1_pismo01_emacron} (\Ju{quidem});
  il. \vref{fig:Ungler1_pismo01_A} (\Ju{Additio}),
  il. \vref{fig:Ungler1_pismo01_D} (\Ju{Deinde}),
  il. \vref{fig:Ungler1_pismo01_M} (\Ju{Mediatio}),
  il. \vref{fig:Ungler1_pismo01_N} (\Ju{Notandum}),
  il. \vref{fig:Ungler1_pismo01_R} (\Ju{Radix}, \Ju{Radicem}),
  il. \vref{fig:Ungler1_pismo01_U} (\Ju{Undi}),
  il. \vref{fig:Ungler1_pismo01_emacron} (\Ju{quidem}),
il. \vref{fig:Ungler1_pismo01_ij} (\Ju{[ne]cessarii}, \Ju{tercii}) i
  il. \vref{fig:Ungler1_pismo01_con} (\Ju{considerata}).

\item Kolejna pozycja to ,,okragłe \textit{d}'' z trudnym do
  identyfikacji diakrytem (kropka? przecinek?)
  %`ꝺ̇?
  --- patrz przykładowe wystąpienia na
  il. \ref{fig:Ungler1_pismo01_dd}. Diakryt najwyraźniej oznacza
  skrócenie.

  \begin{figure}[ht]
  \centering
  \includegraphics[height=10ex]{img/A1521p8Secundo}
  \includegraphics[height=10ex]{img/A1521p9quod}
  \includegraphics[height=10ex]{img/A1521p11secundae}
  \includegraphics[height=10ex]{img/A1521p13de}
  \caption{`{̄ꝺ} z diakrytem: \Ju{Secundo}, \Ju{quod}, \Ju{secundae}?,
    \Ju{de} (\textit{Algorithmus})}
  \label{fig:Ungler1_pismo01_dd}
\end{figure}

  
  \item Litera `e --- patrz przykładowe wystąpienia na
% il. \ref{fig:Ungler1_pismo01_B} (\Ju{Boetius}),
%   il. \ref{fig:Ungler1_pismo01_C} (\Ju{Centum}),
%   il. \ref{fig:Ungler1_pismo01_D} (\Ju{Deinde}),
%   il. \ref{fig:Ungler1_pismo01_D} (\Ju{Decima}),
%   il. \ref{fig:Ungler1_pismo01_D} (\Ju{Differrntia} [sic!]),
%   il. \ref{fig:Ungler1_pismo01_I} (\Ju{Iſte}),
%   il. \ref{fig:Ungler1_pismo01_I} (\Ju{Iſte}),
%   il. \ref{fig:Ungler1_pismo01_M} (\Ju{Impreſſus}),
%   il. \ref{fig:Ungler1_pismo01_M} (\Ju{Mediatio}),
%   %il. \ref{fig:Ungler1_pismo01_M} (\Ju{Materialiter}),
%   il. \ref{fig:Ungler1_pismo01_N} (\Ju{Nec}, \Ju{Numerum}, \Ju{Numeratio}),
%   il. \ref{fig:Ungler1_pismo01_P} (\Ju{Progreſſio}),
%   il. \ref{fig:Ungler1_pismo01_S} (\Ju{Sexto}),
%   il. \ref{fig:Ungler1_pismo01_T} (\Ju{Tres}?, \Ju{Tercia}, \Ju{Tertio}),
%   il. \ref{fig:Ungler1_pismo01_U} (\Ju{Uel}),
%   il. \ref{fig:Ungler1_pismo01_emacron} (\Ju{inuentorem}),
%   il. \ref{fig:Ungler1_pismo01_f} (\Ju{fieri}),
%   il. \ref{fig:Ungler1_pismo01_h} (\Ju{theca}),
il. \vref{fig:Ungler1_pismo01_B} (\Ju{Boetius}),
  il. \vref{fig:Ungler1_pismo01_C} (\Ju{Centum}),
  il. \vref{fig:Ungler1_pismo01_D} (\Ju{Deinde},
\Ju{Decima},
\Ju{Differrntia} [sic!]),
  il. \vref{fig:Ungler1_pismo01_I} (\Ju{Iſte},
\Ju{Iſte}),
  il. \vref{fig:Ungler1_pismo01_M} (\Ju{Impreſſus},
\Ju{Mediatio}),
  %il. \vref{fig:Ungler1_pismo01_M} (\Ju{Materialiter}),
  il. \vref{fig:Ungler1_pismo01_N} (\Ju{Nec}, \Ju{Numerum}, \Ju{Numeratio}),
  il. \vref{fig:Ungler1_pismo01_P} (\Ju{Progreſſio}),
  il. \vref{fig:Ungler1_pismo01_S} (\Ju{Sexto}),
  il. \vref{fig:Ungler1_pismo01_T} (\Ju{Tres}?, \Ju{Tercia}, \Ju{Tertio}),
  il. \vref{fig:Ungler1_pismo01_U} (\Ju{Uel}),
 il. \vref{fig:Ungler1_pismo01_dd} (\Ju{secundae}?),
  il. \vref{fig:Ungler1_pismo01_emacron} (\Ju{inuentorem}),
  il. \vref{fig:Ungler1_pismo01_f} (\Ju{fieri}),
  il. \vref{fig:Ungler1_pismo01_h} (\Ju{theca}),
il. \vref{fig:Ungler1_pismo01_i_} (\Ju{īter}),
il. \vref{fig:Ungler1_pismo01_ij} (\Ju{[ne]cessarii}, \Ju{tercii}),
    il. \vref{fig:Ungler1_pismo01_mmacron} (\Ju{enim}),
   il. \vref{fig:Ungler1_pismo01_omacron} (\Ju{econuerso}),
il. \vref{fig:Ungler1_pismo01_p_} (\Ju{superficies}),
il. \vref{fig:Ungler1_pismo01_pf} (\Ju{proprie}),
  il. \vref{fig:Ungler1_pismo01_con} (\Ju{contentus}, \Ju{completi}, \Ju{considerata}) i
 il. \vref{fig:Ungler1_pismo01_rum} (\Ju{linearum}).


  
\item Litera \textit{e} z kreską `{ē} --- patrz przykładowe wystąpienia na
  il. \vref{fig:Ungler1_pismo01_emacron},
  il. \vref{fig:Ungler1_pismo01_R} (\Ju{Radicem}),
  il. \vref{fig:Ungler1_pismo01_h} (\Ju{habent}),
il. \vref{fig:Ungler1_pismo01_p_} (\Ju{superiorem}) i
il. \vref{fig:Ungler1_pismo01_qet} (\Ju{quocienscumque});
 kreskę nad literą interpretuję jako
  makron, choć w każdym przykładzie ma nieco inna długość (w
  standardzie Unicode \Ju{◌̄} \ucode{0304} \uname{COMBINING MACRON});
  oznacza ona skrócenie.
  \begin{figure}[ht]
  \centering
  \includegraphics[height=10ex]{img/A1521p9pertractionem}
  \includegraphics[height=10ex]{img/A1521p8quidem}
  \includegraphics[height=10ex]{img/A1521p8Inuentorem}
%  \includegraphics[height=10ex]{img/A1521p19Un_}
  \caption{`{ē}: \Ju{pertactionem}, \Ju{quidem}, \Ju{inuentorem}
    (\textit{Algorithmus})}
  \label{fig:Ungler1_pismo01_emacron}
\end{figure}
% p z kreską, t, r,  i, o, n, q, u,  e, rr

\item Litera `f --- patrz przykładowe wystąpienia na
  il. \vref{fig:Ungler1_pismo01_f} i il. \vref{fig:Ungler1_pismo01_p_}
  (\Ju{superiorem},\Ju{superficies}).
\begin{figure}[ht]
  \centering
  \includegraphics[height=10ex]{img/A1521p10fieri}
  \includegraphics[height=10ex]{img/A1521p9figuras}
  \includegraphics[height=10ex]{img/A1521p8formata}
%  \includegraphics[height=10ex]{img/A1521p19Un_}
  \caption{`f: \Ju{fieri}, \Ju{figuras}, \Ju{formata}
    (\textit{Algorithmus})}
  \label{fig:Ungler1_pismo01_f}
\end{figure}

\item Ligatura \textit{ff} --- patrz przykładowe wystąpienie na
  il. \vref{fig:Ungler1_pismo01_D} (\Ju{Differrntia} [sic!]), znak
  występuje w tekście chyba tylko w formach słowa \Ju{differencja} (w
  standardzie Unicode \Ju{ff} \ucode{FB00} \uname{LATIN SMALL LIGATURE
    FF}).
  
\item Litera `g  --- patrz przykładowe wystąpienie na
  il. \vref{fig:Ungler1_pismo01_D} (\Ju{Digitus}),
  il. \vref{fig:Ungler1_pismo01_I} (\Ju{Igitur}),
  il. \vref{fig:Ungler1_pismo01_P} (\Ju{Progessio}) i
  il. \vref{fig:Ungler1_pismo01_f} (\Ju{figuras}).
  
\item Litera `h --- patrz przykładowe wystąpienia na
  il. \vref{fig:Ungler1_pismo01_h}.
\begin{figure}[ht]
  \centering
  \includegraphics[height=10ex]{img/A1521p8habent}
%  \includegraphics[height=10ex]{img/A1521p8nihil}
  \includegraphics[height=10ex]{img/A1521p9ni-hil}
  \includegraphics[height=10ex]{img/A1521p9theca}
%  \includegraphics[height=10ex]{img/A1521p19Un_}
  \caption{`h: \Ju{habent}, \Ju{nihil}, \Ju{theca} [theta] ---
      por. \autocite[s. 3]{sacrobosco39:_johan_sacro_bosco_tract_de}
    (\textit{Algorithmus})}
  \label{fig:Ungler1_pismo01_h}
\end{figure}

\item Litera \textit{h} z trudnym do zidentyfikowania diakrytem
  (przecinek?) --- być może jednak jest to po prostu rozwidlenie
  laski, o którym pisze Porębski \autocite*[s. 11]{porębski05:_paleog}; nie
  znalazłem jej wystapienia w tekście
  \autocite{sacro11:_algor_ioann_de_sacro_busto}, tylko jedno
  wystapienie w \autocite{nieznany10:_rubric_cracov} --- patrz
  il. \vref{fig:Ungler1_pismo01_Rh}. Wiadomo (patrz np.  tabele znaków
  na s. 53 i 62 \autocite{perrousseaux05:_histoir_typog_guten_xviiie}
  i s. 32-33 \autocite{schwenke00:_unter}), że litera ta występowała
  już w biblii Gutenberga, ale nie znam jej funkcji.

   \begin{figure}[ht]
  \centering
  \includegraphics[height=10ex]{img/Rubricella-ephia}
  \caption{\textit{Rubricella}}
  \label{fig:Ungler1_pismo01_Rh} 
\end{figure}

  
\item Litera \textit{i} bez kropki (w standardzie
  Unicode \Ju{ı} \ucode{0131} \uname{LATIN SMALL LETTER DOTLESS
    I}) --- patrz przykładowe wystąpienia
  na il. \vref{fig:Ungler1_pismo01_i-},
  il. \vref{fig:Ungler1_pismo01_D} (\Ju{Decima}),
  il. \vref{fig:Ungler1_pismo01_h} (\Ju{nihil}) i
  il. \vref{fig:Ungler1_pismo01_qet} (\Ju{quocienscumque)}.
  Nie widać jakiejś reguły w jej stosowaniu.

\begin{figure}[ht]
  \centering
  \includegraphics[height=10ex]{img/A1521p9quili-bet}
  \includegraphics[height=10ex]{img/A1521p9i-n}
%  \includegraphics[height=10ex]{img/A1521p8quidam}
%  \includegraphics[height=10ex]{img/A1521p19Un_}
  \caption{`ı: \Ju{quilibet}, \Ju{in}
    (\textit{Algorithmus})}
  \label{fig:Ungler1_pismo01_i-}
\end{figure}

  
\item Litera `i --- patrz przykładowe wystąpienia na
  il. \vref{fig:Ungler1_pismo01_i_}, 
  il. \vref{fig:Ungler1_pismo01_f} (\Ju{fieri}, \Ju{figuras})
    il. \vref{fig:Ungler1_pismo01_h} (\Ju{nihil})
    il. \vref{fig:Ungler1_pismo01_mmacron} (\Ju{omnis}, \Ju{omnibus}),
il. \vref{fig:Ungler1_pismo01_p_} (\Ju{superiorem},\Ju{superficies}),
il. \vref{fig:Ungler1_pismo01_pf} (\Ju{proprie}, \Ju{proximum}),
il. \vref{fig:Ungler1_pismo01_qet} (\Ju{quinque}, \Ju{quocienscumque)},
 il. \vref{fig:Ungler1_pismo01_qetd} (\Ju{quamuis}),
  il. \vref{fig:Ungler1_pismo01_con} \Ju{completi} i
 il. \vref{fig:Ungler1_pismo01_rum} (\Ju{linearum}).

    
  \item Litera \textit{i} bez kropki z kreską `{ı̄} --- patrz
    przykładowe wystąpienia na il. \vref{fig:Ungler1_pismo01_i_},
    kreskę nad literą interpretuję jako makron (znak standardu  Unicode);
%    \Ju{◌̄} \ucode{0304} \uname{COMBINING MACRON});
    oznacza ona
    skrócenie.

\begin{figure}[ht]
  \centering
  \includegraphics[height=10ex]{img/A1521p9i_n}
  \includegraphics[height=10ex]{img/A1521p8i_nter}
%  \includegraphics[height=10ex]{img/A1521p8quidam}
%  \includegraphics[height=10ex]{img/A1521p19Un_}
  \caption{`{ı̄}: \Ju{in}, \Ju{inter}?
    (\textit{Algorithmus})}
  \label{fig:Ungler1_pismo01_i_}
\end{figure}



\item Kolejny niewyraźnie odbity znak interpretuję jako ligaturę
  \textit{i} z \textit{j} pomimo braku drugiej kropki (w standardzie
  Unicode \Ju{ij} \ucode{0133} \uname{LATIN SMALL LIGATURE IJ})---
  patrz przykładowe wystąpienia na il. \vref{fig:Ungler1_pismo01_ij} i
  il. \vref{fig:Ungler1_pismo01_H} (\Ju{His}?).
% także JuniusX[???]}), o co mi chodziło???
\begin{figure}[ht]
  \centering
  \includegraphics[height=10ex]{img/A1521p14-cessari}
  \includegraphics[height=10ex]{img/A1521p14tercij}
  \includegraphics[height=10ex]{img/A1521p14aliis}
%  \includegraphics[height=10ex]{img/A1521p19Un_}
  \caption{`{ı̄}: \Ju{[ne]cessarii}, \Ju{tercii}, \Ju{aliis}
    (\textit{Algorithmus})}
  \label{fig:Ungler1_pismo01_ij}
\end{figure}

% a
% l


\item Litera `k --- nie znalazłem żadnego wystąpienia tego znaku.
  
\item Litera `l --- patrz przykładowe wystąpienia na
  il. \vref{fig:Ungler1_pismo01_M} (\Ju{Moraliter}),
  il. \vref{fig:Ungler1_pismo01_amacron} (\Ju{quamlibet}, \Ju{illam}),
  il. \vref{fig:Ungler1_pismo01_h} (\Ju{nihil}),
  il. \vref{fig:Ungler1_pismo01_i-} (\Ju{quilibet}),
  il. \vref{fig:Ungler1_pismo01_ij} (\Ju{aliis}),
  il. \vref{fig:Ungler1_pismo01_con} \Ju{completi} i
 il. \vref{fig:Ungler1_pismo01_rum} (\Ju{locorum}, \Ju{linearum}).


  
\item Kolejny znak to litera \textit{l} zmodyfikowana (dodatkowy
  kształt może być interpretowany jako poprzeczka lub diakryt
  stykający się z literą) --- nie znalazłem żadnego wystąpienia tego
  znaku w \autocite{sacro11:_algor_ioann_de_sacro_busto}, tylko jedno
  wystapienie w \autocite{nieznany10:_rubric_cracov} --- patrz
 il. \vref{fig:Ungler1_pismo01_Rl}.
 Skądinąd wiadomo, że taki lub podobny znak był używany
  m.in. jako skrót \textit{el}, \textit{ul} lub \textit{vel} --- patrz
  np. \autocite[s. 6]{N3027} i \autocite{bień20:_trakt_stanis_zabor}.
  % ; `ꝉ (znak interpretuję jako
  % wariant graficzny \textit{l} z wysoką
  % poprzeczką\footnote{\ucode{A749} \uname{LATIN SMALL LETTER L WITH
  %     HIGH STROKE}} --- por. Zaborowski),

  \begin{figure}[ht]
  \centering
  \includegraphics[height=10ex]{img/Rubricella-eccX}
  \caption{\textit{Rubricella}}
  \label{fig:Ungler1_pismo01_Rl} 
\end{figure}

% il. \vref{fig:Ungler1_pismo01_Rl}.

  
\item Litera `m --- patrz przykładowe wystąpienia na
    il. \vref{fig:Ungler1_pismo01_C} (\Ju{Centum}, \Ju{Compositus}),
    il. \vref{fig:Ungler1_pismo01_D} (\Ju{Decima}),
    il. \vref{fig:Ungler1_pismo01_I} (\Ju{Impressus}),
    il. \vref{fig:Ungler1_pismo01_N} (\Ju{Numerum}, \Ju{Numeratio}),
    il. \vref{fig:Ungler1_pismo01_O} (\Ju{Omnis}),
    il. \vref{fig:Ungler1_pismo01_P} (\Ju{Primo}),
    il. \vref{fig:Ungler1_pismo01_f} (\Ju{formata}) i
il. \vref{fig:Ungler1_pismo01_pf} (\Ju{proximum}).

\item Litera \textit{m} z kreską --- patrz przykładowe wystąpienia na
    il. \vref{fig:Ungler1_pismo01_mmacron} i
    il. \vref{fig:Ungler1_pismo01_O} (\Ju{Omnis});
%    
    kreskę nad literą interpretuję jako  makron standardzie, choć czasami ma nieco inna długość;
  %   (w
  % standardzie Unicode \Ju{◌̄} \ucode{0304} \uname{COMBINING MACRON});
  oznacza ona skrócenie.

  \begin{figure}[ht]
  \centering
  \includegraphics[height=10ex]{img/A1521p8omnis}
  \includegraphics[height=10ex]{img/A1521p16tomnibus}
  \includegraphics[height=10ex]{img/A1521p8em_}
  \includegraphics[height=10ex]{img/A1521p14-prima-figura}
  \caption{`{m̄}: \Ju{omnis}, \Ju{omnibus}, \Ju{enim}? \Ju{Q?? [sia prima figura]}
    (\textit{Algorithmus})}
  \label{fig:Ungler1_pismo01_mmacron} 
\end{figure}

% o, i, b, u, s, e, Q

\item Litera `n --- patrz przykładowe wystąpienia na
   il. \vref{fig:Ungler1_pismo01_C} (\Ju{Centum}),
   il. \vref{fig:Ungler1_pismo01_D} (\Ju{Deinde}, \Ju{Differrntia}[sic]),
   il. \vref{fig:Ungler1_pismo01_H} (\Ju{Hanc}),
   il. \vref{fig:Ungler1_pismo01_M} (\Ju{Minus}),
   il. \vref{fig:Ungler1_pismo01_N} (\Ju{Notandum}),
   il. \vref{fig:Ungler1_pismo01_Q} (\Ju{Quinta}),
   il. \vref{fig:Ungler1_pismo01_S} (\Ju{Sinistrorsum}),
   il. \vref{fig:Ungler1_pismo01_emacron} (\Ju{inuentorem}),
   il. \vref{fig:Ungler1_pismo01_h} (\Ju{nihil}),
   il. \vref{fig:Ungler1_pismo01_i-} (\Ju{in}),
   il. \vref{fig:Ungler1_pismo01_omacron} (\Ju{non}),
  il. \vref{fig:Ungler1_pismo01_con} (\Ju{contentus}) i
 il. \vref{fig:Ungler1_pismo01_rum} (\Ju{linearum}).


 \item Litera \textit{n} z kreską --- patrz przykładowe wystąpienia na
   il. \vref{fig:Ungler1_pismo01_nmacron} i
   il. \vref{fig:Ungler1_pismo01_Q} (\Ju{Quinam});  kreskę nad literą
   interpretuję jako makron (znak standardu Unicode);
   % \Ju{◌̄} \ucode{0304}
   % \uname{COMBINING MACRON});
   oznacza ona skrócenie.

  \begin{figure}[ht]
  \centering
  \includegraphics[height=10ex]{img/A1521p10tamen}
  \includegraphics[height=10ex]{img/A1521p23possunt}
  % \includegraphics[height=10ex]{img/A1521p8em_}
  % \includegraphics[height=10ex]{img/A1521p14-prima-figura}
  \caption{`{n̄}: \Ju{tamen}?, \Ju{possunt} (\textit{Algorithmus})}
  \label{fig:Ungler1_pismo01_nmacron} 
\end{figure}

% p
% t

\item Kolejny znak to litera \textit{n} z dwiema kropkami --- nie znalazłem żadnego wystąpienia tego
  znaku
w tekście
  \autocite{sacro11:_algor_ioann_de_sacro_busto}, tylko  wystapienia
  w \autocite{nieznany10:_rubric_cracov} --- patrz
  il. \vref{fig:Ungler1_pismo01_Rn}.  % `{n̈} (diereza),

     \begin{figure}[ht]
  \centering
  \includegraphics[height=10ex]{img/Rubricella-dnica}
  \caption{\textit{Rubricella}}
  \label{fig:Ungler1_pismo01_Rn} 
\end{figure}


  
\item Litera `o --- patrz przykładowe wystąpienia na
   il. \vref{fig:Ungler1_pismo01_mmacron} (\Ju{omnis}, \Ju{omnibus}),
   il. \vref{fig:Ungler1_pismo01_omacron} (\Ju{econuerso}),
il. \vref{fig:Ungler1_pismo01_p_} (\Ju{superiorem}),
il. \vref{fig:Ungler1_pismo01_pf} (\Ju{propositus}),
il. \vref{fig:Ungler1_pismo01_qet} (\Ju{quocienscumque}) i
 il. \vref{fig:Ungler1_pismo01_rum} (\Ju{locorum}).


\item Litera \textit{o} z kreską --- patrz przykładowe wystąpienia na
   il. \vref{fig:Ungler1_pismo01_omacron};
%
  kreskę nad literą interpretuję jako makron (znak standardu Unicode);
  %\Ju{◌̄} \ucode{0304} \uname{COMBINING MACRON});
  oznacza ona
  skrócenie.

  \begin{figure}[ht]
  \centering
  \includegraphics[height=10ex]{img/A1521p20non}
  \includegraphics[height=10ex]{img/A1521p20potest}
  \includegraphics[height=10ex]{img/A1521p20econuerso}
  % \includegraphics[height=10ex]{img/A1521p14-prima-figura}
  \caption{`{ō}: \Ju{non}, \Ju{potest}, \Ju{econuerso}
    (\textit{Algorithmus})}
  \label{fig:Ungler1_pismo01_omacron} 
\end{figure}

% n p t e c u r logs o


\item Litera `p --- patrz przykładowe wystąpienia na
   il. \vref{fig:Ungler1_pismo01_nmacron} (\Ju{possunt}),
   il. \vref{fig:Ungler1_pismo01_omacron} (\Ju{potest}),
il. \vref{fig:Ungler1_pismo01_pf} (\Ju{proprie},\Ju{propositus}) i
  il. \vref{fig:Ungler1_pismo01_con} (\Ju{completi}).

 \end{enumerate}

\item Wiersz drugi (w wykazie trzeci):
  \begin{enumerate}
  \item Litera \textit{p} z kropką
  --- nie znalazłem jej wystąpienia w tekście.  
%`{ṗ} (\textit{p} z kropką), 

\item Kolejny znak to niewyraźnie odbity brewigraf \textit{p} z
  przekreślonym wydłużeniem dolnym (w standardzie Unicode \ucode{A751}
  \uname{LATIN SMALL LETTER P WITH STROKE THROUGH DESCENDER}) ---
  patrz przykładowe wystąpienia na il. \vref{fig:Ungler1_pismo01_p_};
  we wszystkich użyciach oznacza \textit{per}.

\begin{figure}[ht]
  \centering
  \includegraphics[height=10ex]{img/A1521p19per}
  \includegraphics[height=10ex]{img/A1521p17superiorem}
  \includegraphics[height=10ex]{img/A1521p20superficies}
%  \includegraphics[height=10ex]{img/A1521p19Un_}
  \caption{`ꝑ: \Ju{per}, \Ju{superiorem},\Ju{superficies}
    (\textit{Algorithmus})}
  \label{fig:Ungler1_pismo01_p_}
\end{figure}

%long s
% u
% i
% o
% r
% e
% e _
% f
% c


\item Kolejny znak to brewigraf, który nazywam literą \textit{p} z
  floresem (por. %uzasadnienie przedstawiłem w
  \autocite{bień20:_trakt_stanis_zabor}, w standardzie Unicode `ꝓ
  \ucode{A753} \uname{LATIN SMALL LETTER P WITH FLOURISH}) --- patrz
  przykładowe wystąpienia na na il. \vref{fig:Ungler1_pismo01_pf}, we
  wszystkich użyciach oznacza \textit{pro}.
  
 \begin{figure}[ht]
  \centering
  \includegraphics[height=10ex]{img/A1521p21proprie}
  \includegraphics[height=10ex]{img/A1521p22propositus}
  \includegraphics[height=10ex]{img/A1521p19proximum}
%  \includegraphics[height=10ex]{img/A1521p19Un_}
  \caption{`ꝓ:  \Ju{proprie},\Ju{propositus}, \Ju{proximum}
    (\textit{Algorithmus})}
   \label{fig:Ungler1_pismo01_pf}
 \end{figure}

% p r i e o ligatura si m x t us uz

  \item Litera `q --- patrz przykładowe wystąpienia na
  il. \vref{fig:Ungler1_pismo01_amacron} (\Ju{quamlibet}, \Ju{quidam}),
  il. \vref{fig:Ungler1_pismo01_dd} (\Ju{quod}),
  il. \vref{fig:Ungler1_pismo01_emacron} (\Ju{quidem}),
  il. \vref{fig:Ungler1_pismo01_i-} (\Ju{quilibet}) i
il. \vref{fig:Ungler1_pismo01_qet} (\Ju{quinque}, \Ju{quocienscumque}).


\item Kolejna pozycja to litera \textit{q} z kreską --- patrz 
  znalezione wystąpienia na il. \vref{fig:Ungler1_pismo01_q_},
%
  kreskę nad literą interpretuję jako makron (znak standard Unicode);
  %\Ju{◌̄} \ucode{0304} \uname{COMBINING MACRON});
  oznacza ona
  skrócenie.

   \begin{figure}[ht]
  \centering
  \includegraphics[height=10ex]{img/A1521p14quia}
 \includegraphics[height=10ex]{img/A1521p9quae}
  % \includegraphics[height=10ex]{img/A1521p19proximum}
%  \includegraphics[height=10ex]{img/A1521p19Un_}
  \caption{`{q̄}:  \Ju{quia}?, \Ju{quae}?
    (\textit{Algorithmus})}
   \label{fig:Ungler1_pismo01_q_}
 \end{figure}

%`{q̄} (macron)

\item Kolejny znak to brewigraf w postaci litery \textit{q}
  z wydłużeniem dolnym przekreślonym ukośną kreską u góry zagiętą w lewo (w standardzie Unicode
  pomimo nieco innego kształtu w niektórych fontach znak można
  reprezentować jako `ꝙ \ucode{A759} \uname{LATIN SMALL LETTER Q WITH
    DIAGONAL STROKE})--- patrz
  przykładowe wystąpienia na na il. \vref{fig:Ungler1_pismo01_qku}, we
  wszystkich użyciach oznacza \textit{quod} .
  
   \begin{figure}[ht]
  \centering
  \includegraphics[height=10ex]{img/A1521p21quod}
  \includegraphics[height=10ex]{img/A1521p23quod}
%  \includegraphics[height=10ex]{img/A1521p19proximum}
%  \includegraphics[height=10ex]{img/A1521p19Un_}
  \caption{`ꝙ:  \Ju{quod}
    (\textit{Algorithmus})}
   \label{fig:Ungler1_pismo01_qku}
 \end{figure}

    
\item Kolejny znak to brewigraf w formie ligatury litery \textit{q} i
  znaku, który interpetuję jako formę brewigrafu \textit{et}
(w standardzie Unicode znak można reprezentować za
  pomocą tzw. znaku prywatnego zgodnego z rekomendacja MUFI
  \autocite[s. 81]{haugen15:_mufi1} ` \mcode{F8BF} \mname{LATIN SMALL LETTER Q
    LIGATED WITH FINAL ET})  --- patrz przykładowe wystąpienia na na
  il. \vref{fig:Ungler1_pismo01_qet}, we wszystkich użyciach oznacza
  \textit{que}. 

     \begin{figure}[ht]
  \centering
  \includegraphics[height=10ex]{img/A1521p9usque}
  \includegraphics[height=10ex]{img/A1521p25quinque}
  \includegraphics[height=10ex]{img/A1521p24quociescumque}
%  \includegraphics[height=10ex]{img/A1521p19Un_}
  \caption{`:  \Ju{vsque}, \Ju{quinque}, \Ju{quocienscumque}
    (\textit{Algorithmus})}
   \label{fig:Ungler1_pismo01_qet}
 \end{figure}

% v long s q u i o c i- e_ u-
  
\item Kolejny znak to też brewigraf w formie ligatury litery
  \textit{q} i znaku, który interpetuję jako formę brewigrafu
  \textit{et}, ale z trudnym do identyfikacji diakrytem --- patrz
  przykładowe wystąpienia na na il. \vref{fig:Ungler1_pismo01_qetd},
  oznacza \textit{quan} lub \textit{quam}. Ze względu na diakryt
  reprezentacja w standardzie Unicode nie jest oczywista, jedna z
  możliwości to interpretowanie diakrytu jako nadpisanego spłaszczonego
  otwartego \textit{a}: \ucode{1DD3} \uname{COMBINING LATIN SMALL LETTER
    FLATTENED OPEN A ABOVE}: `{ᷓ} --- por.
  \autocite{bień20:_trakt_stanis_zabor}.

     \begin{figure}[ht]
  \centering
  \includegraphics[height=10ex]{img/A1521p21quantum}
  \includegraphics[height=10ex]{img/A1521p9quamuis}
  \includegraphics[height=10ex]{img/A1521p23quam} 
%  \includegraphics[height=10ex]{img/A1521p19Un_}
  \caption{\Ju{quantum}, \Ju{quamuis}, \Ju{quam}
    (\textit{Algorithmus})}
   \label{fig:Ungler1_pismo01_qetd}
 \end{figure}

% t u_ u i s


\item Litera `r --- patrz przykładowe wystąpienia na
    il. \vref{fig:Ungler1_pismo01_A} (\Ju{Articulus}),
    il. \vref{fig:Ungler1_pismo01_D} (\Ju{Differrntia} [sic]),
    il. \vref{fig:Ungler1_pismo01_I} (\Ju{Igitur}),
    il. \vref{fig:Ungler1_pismo01_M} (\Ju{Materialiter}),
    il. \vref{fig:Ungler1_pismo01_N} (\Ju{Numerum}, \Ju{Numeratio}),
    il. \vref{fig:Ungler1_pismo01_P} (\Ju{Progressio}),
    il. \vref{fig:Ungler1_pismo01_Q} (\Ju{Quarta}),
    il. \vref{fig:Ungler1_pismo01_S} (\Ju{Subtratio} [sic!], \Ju{Similiter}),
    il. \vref{fig:Ungler1_pismo01_emacron} (\Ju{pertactionem}),
    il. \vref{fig:Ungler1_pismo01_f} (\Ju{fieri}, \Ju{figuras}), 
    il. \vref{fig:Ungler1_pismo01_i_} (\Ju{inter}),
il. \vref{fig:Ungler1_pismo01_ij} (\Ju{[ne]cessarii}, \Ju{tercii}),
   il. \vref{fig:Ungler1_pismo01_omacron} (\Ju{econuerso}) i
  il. \vref{fig:Ungler1_pismo01_con} (\Ju{considerata}).



\item Kolejny znak to ,,okrągłe \textit{r}'' (w standardzie Unicode
  \`ꝛ  \ucode{A75B} \uname{LATIN SMALL LETTER R ROTUNDA})
%  (\textit{r rotunda})
--- patrz przykładowe wystąpienia na
    il. \vref{fig:Ungler1_pismo01_P} (\Ju{Primo}),
    il. \vref{fig:Ungler1_pismo01_S} (\Ju{Sinistrorsum}),
    il. \vref{fig:Ungler1_pismo01_emacron} (\Ju{inuentorem}),
    il. \vref{fig:Ungler1_pismo01_p_} (\Ju{superiorem}) i
    il. \vref{fig:Ungler1_pismo01_pf} (\Ju{proprie}).

\item  Litera `s: --- patrz przykładowe wystąpienia na
  il. \vref{fig:Ungler1_pismo01_L} (\Ju{Locus}),
  il. \vref{fig:Ungler1_pismo01_M} (\Ju{Minus}),
  il. \vref{fig:Ungler1_pismo01_O} (\Ju{Omnis}),
  il. \vref{fig:Ungler1_pismo01_T} (\Ju{Tres}),
  il. \vref{fig:Ungler1_pismo01_U} (\Ju{Unitas}),
  il. \vref{fig:Ungler1_pismo01_f} (\Ju{figuras}),
  il. \vref{fig:Ungler1_pismo01_ij} (\Ju{aliis}),
      il. \vref{fig:Ungler1_pismo01_mmacron} (\Ju{omnibus}),
      il. \vref{fig:Ungler1_pismo01_p_} (\Ju{superficies}),
 il. \vref{fig:Ungler1_pismo01_qetd} (\Ju{quamuis}) i
  il. \vref{fig:Ungler1_pismo01_con} (\Ju{contentus}).

  
\item Kolejna pozycja to litera \textit{s} z kreską--- nie znalazłem
  jej wystąpienia w tekście.
% `{s̄} (makron), 

\item Litera długie \textit{s} (w standardzie Unicode `ſ \ucode{017F}
  \uname{LATIN SMALL LETTER LONG S}) --- patrz przykładowe wystąpienia
  na il. \vref{fig:Ungler1_pismo01_C} (\Ju{Compositus}),
  il. \vref{fig:Ungler1_pismo01_S} (\Ju{Sinistrorsum}),
  il. \vref{fig:Ungler1_pismo01_dd} (\Ju{secundae}?),
  il. \vref{fig:Ungler1_pismo01_omacron} (\Ju{econuerso}) i
  il. \vref{fig:Ungler1_pismo01_qet} (\Ju{vsque},
  \Ju{quocienscumque}).


\item Kolejna pozycja to ligatura długiego \textit{s} ze znakiem,
  który jest chyba brewigrafem \textit{et} --- patrz przykładowe
  wystąpienia na il. \ref{fig:Ungler1_pismo01_sed}, chyba zawsze
  oznacza \textit{sed}.


% (w standardzie Unicode znak można reprezentować za
%   pomocą tzw. znaku prywatnego zgodnego z rekomendacja MUFI
%   \autocite[s. 81]{haugen15:_mufi1} ` \mcode{F8BF} \mname{LATIN SMALL LETTER Q
%     LIGATED WITH FINAL ET}). 
%   `ſ (długie \textit{s}) `ꝫ (czyli literę
% \textit{et}\footnote{\ucode{A76B} \uname{LATIN SMALL LETTER ET}})

     \begin{figure}[ht]
  \centering
  \includegraphics[height=10ex]{img/A1521p12sed1}
  \includegraphics[height=10ex]{img/A1521p10sed}
  \includegraphics[height=10ex]{img/A1521p12sed}
%  \includegraphics[height=10ex]{img/A1521p19Un_}
  \caption{\Ju{sed}
    (\textit{Algorithmus})}
   \label{fig:Ungler1_pismo01_sed}
 \end{figure}


\item Kolejna pozycja to ligatura długiego \textit{s} ze znakiem
  omawianym wcześniej, mianowicie literą \textit{b} z
  niezidentyfikowanym diakrytem --- patrz przykładowe wystąpienia na
  il. \ref{fig:Ungler1_pismo01_sub}, chyba zawsze oznacza
  \textit{sub}.

% `ſ (długie \textit{s}) chyba występująca wcześniej samodzielnie
% litera \textit{b} z diakrytem.

     \begin{figure}[ht]
  \centering
  \includegraphics[height=10ex]{img/A1521p22sub1}
  \includegraphics[height=10ex]{img/A1521p22sub}
  \includegraphics[height=10ex]{img/A1521p23sub}
%  \includegraphics[height=10ex]{img/A1521p19Un_}
  \caption{\Ju{sub}
    (\textit{Algorithmus})}
  \label{fig:Ungler1_pismo01_sub}
 \end{figure}


\item Kolejna pozycja to ligatura dwóch długich \textit{s} (w
  standardzie Unicode znak można reprezentować za pomocą tzw. znaku
  prywatnego zgodnego z rekomendacja MUFI
  \autocite[s. 88]{haugen15:_mufi1} ` \mcode{EBA6} \mname{LATIN SMALL
    LIGATURE LONG S LONG S}) --- patrz przykładowe wystąpienia na na
  il. \vref{fig:Ungler1_pismo01_I} (\Ju{Impressus}),
  il. \vref{fig:Ungler1_pismo01_P} (\Ju{Progreſſio}) i
  il. \vref{fig:Ungler1_pismo01_ij} (\Ju{[ne]cessarii}) .

\item Kolejna pozycja to ligatura \textit{s} i \textit{t} (w standardzie Unicode
`ſt \ucode{FB05} \uname{LATIN SMALL LIGATURE LONG S T})
--- patrz przykładowe wystąpienie na
il. \vref{fig:Ungler1_pismo01_S} (\Ju{Sinistrorsum}).


\item Litera `t --- patrz przykładowe wystąpienia na
    il. \vref{fig:Ungler1_pismo01_A} (\Ju{Articulus}),
    il. \vref{fig:Ungler1_pismo01_B} (\Ju{Boetius}),
    il. \vref{fig:Ungler1_pismo01_C} (\Ju{Centum}),
      il. \vref{fig:Ungler1_pismo01_D} (\Ju{Digitus}),
      il. \vref{fig:Ungler1_pismo01_E} (\Ju{Et}),
        il. \vref{fig:Ungler1_pismo01_I} (\Ju{Igitur}),
    il. \vref{fig:Ungler1_pismo01_N} (\Ju{Notandum}, \Ju{Numeratio}),
    il. \vref{fig:Ungler1_pismo01_Q} (\Ju{Quarta}, \Ju{Quinta}),
    il. \vref{fig:Ungler1_pismo01_S} (\Ju{Sexto}, \Ju{Subtratio} [sic!]),
  il. \vref{fig:Ungler1_pismo01_U} (\Ju{Unitas}),
  il. \vref{fig:Ungler1_pismo01_emacron} (\Ju{pertactionem}),
  il. \vref{fig:Ungler1_pismo01_f} (\Ju{formata}),
  il. \vref{fig:Ungler1_pismo01_h} (\Ju{habent}, \Ju{theca}),
  il. \vref{fig:Ungler1_pismo01_i-} (\Ju{quilibet}),
il. \vref{fig:Ungler1_pismo01_i_} (\Ju{īter}),
il. \vref{fig:Ungler1_pismo01_ij} (\Ju{tercii}),
   il. \vref{fig:Ungler1_pismo01_nmacron} (\Ju{tamen}?, \Ju{possunt}),
   il. \vref{fig:Ungler1_pismo01_omacron} (\Ju{potest}),
il. \vref{fig:Ungler1_pismo01_pf} (\Ju{propositus}),
 il. \vref{fig:Ungler1_pismo01_qetd} (\Ju{quantum}) i
  il. \vref{fig:Ungler1_pismo01_con} (\Ju{contentus}, \Ju{completi}).

\item Kolejny znak to litera \textit{t} z budzącym wątpliwości
  diakrytem --- nie znalazłem w tekście żadnego wystąpienia, co
  utrudnia interpretację. Nie wykluczam, że diakryt to znak
  zdefiniowany w standardzie Unicode jako \Ju{◌᷑} \ucode{1DD1}
  \uname{COMBINING UR ABOVE}.

\item Litera `u --- patrz przykładowe wystąpienia na
    il. \vref{fig:Ungler1_pismo01_A} (\Ju{Articulus}),
il. \vref{fig:Ungler1_pismo01_B} (\Ju{Boetius}),
  il. \vref{fig:Ungler1_pismo01_C} (\Ju{Centum}),
  il. \vref{fig:Ungler1_pismo01_D} (\Ju{Digitus}),
  il. \vref{fig:Ungler1_pismo01_H} (\Ju{Huius}),
  il. \vref{fig:Ungler1_pismo01_I} (\Ju{Igitur}, \Ju{Impreſſus}),
    il. \vref{fig:Ungler1_pismo01_L} (\Ju{Locus}),
  il. \vref{fig:Ungler1_pismo01_M} (\Ju{Maius}, \Ju{Minus}),
    il. \vref{fig:Ungler1_pismo01_Q} (\Ju{Quinta}),
    il. \vref{fig:Ungler1_pismo01_S} (\Ju{Subtratio} [sic!]),
  il. \vref{fig:Ungler1_pismo01_emacron} (\Ju{quidem}, \Ju{inuentorem}),
  il. \vref{fig:Ungler1_pismo01_f} (\Ju{figuras}),
  il. \vref{fig:Ungler1_pismo01_i-} (\Ju{quilibet}),
   il. \vref{fig:Ungler1_pismo01_omacron} (\Ju{econuerso}),
il. \vref{fig:Ungler1_pismo01_p_} (\Ju{superiorem},\Ju{superficies}),
il. \vref{fig:Ungler1_pismo01_qet} (\Ju{quinque}, \Ju{quocienscumque}),
 il. \vref{fig:Ungler1_pismo01_qetd} (\Ju{quamuis}) i
  il. \vref{fig:Ungler1_pismo01_con} (\Ju{contentus}).

\item Litera \textit{u} z kreską `{ū} --- patrz przykładowe
  wystąpienia na il. \vref{fig:Ungler1_pismo01_C} (\Ju{Cum}),
  il. \vref{fig:Ungler1_pismo01_N} (\Ju{Numerum}, \Ju{Notandum}),
  il. \vref{fig:Ungler1_pismo01_pf} (\Ju{proximum}),
  il. \vref{fig:Ungler1_pismo01_qet} (\Ju{quocienscumque}) i
  il. \vref{fig:Ungler1_pismo01_qetd} (\Ju{quantum}); kreskę nad literą
  interpretuję jako makron (znak standardu standardu Unicode \Ju{◌̄});
  % \ucode{0304}
  % \uname{COMBINING MACRON});
  oznacza ona skrócenie.

 
\item Litera `v  --- patrz przykładowe
  wystąpienie
il. \vref{fig:Ungler1_pismo01_qet} (\Ju{vsque}).

\item Kolejna litera to \textit{v} z kropką lub podobnym diakrytem
%  `{v̇} (kropka?),
---  nie znalazłem
  jej wystąpienia w tekście.

\item Litera `w ---  nie znalazłem
  jej wystąpienia w tekście.

\item Kolejny znak to zapewne minuskuła ,,ogoniaste z''
  (ang. \textit{tailed z}); nie znalazłem jej wystąpienia w tekście. W
  standardzie Unicode od wersji 3.0 jest ona reprezentowana jako `ʒ
  \ucode {0292} \uname{LATIN SMALL LETTER EZH}
  \autocite[113-114]{haugen15:_mufi1}.
  
\item Kolejny znak to brewigraf przypominajacy odwróconą literę
  \textit{c} (w standardzie Unicode `ↄ \ucode{2184} \uname{LATIN SMALL
    LETTER REVERSED C}; rekomendacja MUFI używa też nazwy \mname{LATIN
    ABBREVIATION SIGN SMALL CON} \autocite[s. 29]{haugen15:_mufi1})
  --- patrz przykładowe wystąpienia na
  il. \vref{fig:Ungler1_pismo01_con}, oznacza \textit{con} lub
  \textit{com}.

 \begin{figure}[ht]
  \centering
  \includegraphics[height=10ex]{img/A1521p23contentus}
  \includegraphics[height=10ex]{img/A1521p23completi}
  \includegraphics[height=10ex]{img/A1521p25considera}
%  \includegraphics[height=10ex]{img/A1521p19Un_}
  \caption{\Ju{contentus}, \Ju{completi}, \Ju{considera}
    (\textit{Algorithmus})}
  \label{fig:Ungler1_pismo01_con}
 \end{figure}



% t
%  e n u s p l i d r a
 
 
\item Ostatni znak drugiego wiersza zestawu to niewatpliwie nota
  tyrońska \textit{et} (w standardzie Unicode
  \ucode{204A} \uname{TIRONIAN SIGN ET}, kształt zależy od fontu,
  np. `⁊ ) --- patrz przykładowe wystąpienia
  il. \vref{fig:Ungler1_pismo01_et}. 
   \begin{figure}[ht]
  \centering
  \includegraphics[height=10ex]{img/A1521p25et}
  \includegraphics[height=10ex]{img/A1521p22etc}
%  \includegraphics[height=10ex]{img/A1521p25considera}
%  \includegraphics[height=10ex]{img/A1521p19Un_}
  \caption{\Ju{et}, \Ju{et cetera}
    (\textit{Algorithmus})}
  \label{fig:Ungler1_pismo01_et}
 \end{figure}


 \end{enumerate}
\end{enumerate}



\section{Inne znaki}
\label{sec:inne-znaki}


Czwarty wiersz zestawu zawiera tylko pięć znaków, a piąty i zarazem
ostatni 12 znaków.
\begin{itemize}
\item Wiersz czwarty
  \begin{enumerate}
  \item Pierwszy znak to moim zdaniem brewigraf \textit{rum} (w
    standardzie Unicode \ucode{A75D} \uname{LATIN SMALL LETTER RUM
      ROTUNDA}, kształt zależy od fontu, np. `ꝝ ) --- patrz
    przykładowe wystąpienia na il. \vref{fig:Ungler1_pismo01_rum}. 
   \begin{figure}[ht]
  \centering
  \includegraphics[height=10ex]{img/A1521p25locorum}
  \includegraphics[height=10ex]{img/A1521p25linearum}
%  \includegraphics[height=10ex]{img/A1521p25considera}
%  \includegraphics[height=10ex]{img/A1521p19Un_}
  \caption{\Ju{locorum}, \Ju{linearum}
    (\textit{Algorithmus})}
  \label{fig:Ungler1_pismo01_rum}
 \end{figure}

 % il. \vref{fig:Ungler1_pismo01_rum} \Ju{locorum}, \Ju{linearum}
 % l
 % o
 % c
 % i
 % n
 % e
 % a
\item Następny znak to moim zdaniem brewigraf \textit{us} (w standardzie Unicode
`ꝰ \ucode{A770} \uname{MODIFIER LETTER US}) --- patrz
  przykładowe wystąpienia na
il. \vref{fig:Ungler1_pismo01_C} (\Ju{Compositus}) i
il. \vref{fig:Ungler1_pismo01_pf} (\Ju{propositus}).

\item Kolejny znak to niewyraźnie odbity podwójny
pochyły dywiz (w standardzie Unicode `⸗
\ucode{2E17} \uname{DOUBLE OBLIQUE HYPHEN})

\item Kolejny znak to dwukropek --- nie znalazłem żadnego wystąpienia tego
  znaku.
  
\item Ostatni znak to kropka.
\end{enumerate}
\item Wiersz piąty (ostatni) zawiera znaki, które nie należą do pisma
  nr 1, ale prawdopodobnie zostały umieszczone w zestawie niejako przy
  okazji.
  \begin{enumerate}
  \item Pierwszy znak to rubryka oznaczana w tym zeszycie za
    Piekarskim grecką literą \textbf{α}.
  \item Następnych 10 pozycji to cyfry zestawu nr 1 --- zestawy cyfr
    są numerowane niezależnie od pism.
  \item Ostatni znak to asteryks, który nie został w żaden sposób
    skomentowany --- nie znalazłem go w tekście.
   \end{enumerate}
\end{itemize}


% COMBINING LATIN SMALL LETTER FLATTENED OPEN A ABOVE is used for ua in qà qua ‘as’, gà rda
% guarda ‘guard’, for ra or ar in Latin contà contra ‘against’, supà supra ‘above’, Portuguese comp
% compra ‘a purchase’, mà ia maria ‘Maria’, pà te parte ‘part’, for numerals và quinta ‘fifth’, ı à prima
% ‘first’, una ‘one’ (Figures 33, 47, 67, 73, 76).

% https://medievalfragments.wordpress.com/2014/05/24/the-punctus-and-his-friends-medieval-punctuation/
% https://sites.ualberta.ca/~sreimer/ms-course/course/punc.htm

% \section{Znaki pominięte}
% \label{sec:znaki-pominite}


\end{document}


%%% Local Variables:
%%% mode: latex
%%% ispell-local-dictionary: "polish"
%%% TeX-master: "Ungler"
%%% coding: utf-8-unix
%%% TeX-PDF-mode: t
%%% TeX-engine: xetex
%%% End:
